% !PRINTED 307
\weberpage{307}

\gesperrt{Thus neither a $VP$ nor a $PP$ is lost or gained.}

But if $f(x)$ itself passes through zero, then, assuming $f_1(x)$ positive, which suffices in view of the remark made above, we have
\[
\begin{array}{c|cc}
 & 0 & 1 \\ \hline
f(x-0) & - & + \\
f(x+0) & + & + \\ \hline
F & + & +
\end{array}
\text{.}
\]
Here, then, a $VP$ passes into a $PP$; that is, \gesperrt{a $VP$ is lost and a $PP$ gained.}

If, finally, $F_\nu$ passes through zero, where of course we must have $0 < \nu < n$, the behaviour is as follows:

If $F_\nu = 0$, then $f_{\nu-1}$ and $f_{\nu+1}$ must necessarily have the same sign, since otherwise $F_\nu$ would consist of two positive terms and could not vanish. Further, by (6), $F_\nu'(x)$ has the same sign as the product $f_{\nu-1} f_\nu F_{\nu+1}$, and hence $F_\nu(x-0)$ has the opposite sign to the product, and $F_\nu(x+0)$ the same sign as it.

Here too we may omit from the enumeration those cases which arise from the ones listed by a simultaneous change of sign in all three functions $F_{\nu-1}, F_\nu, F_{\nu+1}$; there remain, then, the following cases:
% [sic] The first right-hand header is printed with a prime after $\nu+1$; the corresponding later header is unprimed.
\[
\begin{array}{c|ccc||ccc}
f & \nu-1, & \nu, & \nu+1 & \nu-1, & \nu, & \nu+1' \\
 & + & + & + & + & + & + \\ \hline
F(x-0) & + & - & + & - & - & + \\ \hline
F(x+0) & + & + & + & - & + & +
\end{array}
\]
\[
\begin{array}{c|ccc||ccc}
f & \nu-1, & \nu, & \nu+1 & \nu-1, & \nu, & \nu+1 \\
 & + & - & + & + & - & + \\ \hline
F(x-0) & + & + & + & - & + & + \\ \hline
F(x+0) & + & - & + & - & - & +
\end{array}
\text{.}
\]
The changes of sign in these four cases are:
%CONT>>
